\textbf{Abstract}

This report presents the development of the Helike \#213 1P PocketQube
satellite, designed by team Serra Rocketry (IPRJ-UERJ) for the
Satellite Challenge category of the 7th Latin American Space Challenge
(LASC 2026). The mission's central element is the SRAB (Bioinspired
Autorotative Recovery System, from the Portuguese \textit{Sistema de Recuperação
Autorrotativo Bioinspirado}) --- a passive, parachute-free recovery
mechanism based on the autorotation of samara (maple) seeds, which
offers predictable, modelable aerodynamic behavior without the
electromechanical sequencers, pyrotechnics and tangled lines that make
conventional parachutes fail. The satellite also carries a LoRa
communication system with RSSI trilateration from three ground
beacons, which helps locate the satellite after landing and supports
recovery rather than being the mission's main objective. The report
walks through the system architecture and concept of operations, the
physics of samara autorotation, three drop-test campaigns, and the
mathematical modeling and computational simulation of the descent from
the apogee of the team's Dédalo launch vehicle, validated with
RocketPy --- closing with the risk analysis (P $\times$ S matrix) and the pre-flight
validation plan.

\textbf{Authors}

Diego dos Santos Alves. \emph{Serra Rocketry, Instituto Politécnico da UERJ
(IPRJ-UERJ), Nova Friburgo, Rio de Janeiro, 28625-570, Brazil.}

Pedro Henrique Couto Silva. \emph{Serra Rocketry, Instituto Politécnico da
UERJ (IPRJ-UERJ), Nova Friburgo, Rio de Janeiro, 28625-570, Brazil.}

Vinicius Carvalho Monnerat Bandeira. \emph{Serra Rocketry, Instituto
Politécnico da UERJ (IPRJ-UERJ), Nova Friburgo, Rio de Janeiro, 28625-570,
Brazil.}

Caio Phelipe Souza De Lima. \emph{Serra Rocketry, Instituto Politécnico da
UERJ (IPRJ-UERJ), Nova Friburgo, Rio de Janeiro, 28625-570, Brazil.}

João Victor Domingues Barizon de Mello. \emph{Serra Rocketry, Instituto
Politécnico da UERJ (IPRJ-UERJ), Nova Friburgo, Rio de Janeiro, 28625-570,
Brazil.}

Pablo Alves Mattos Borges. \emph{Serra Rocketry, Instituto Politécnico da
UERJ (IPRJ-UERJ), Nova Friburgo, Rio de Janeiro, 28625-570, Brazil.}

Angelo Mondaini Calvão. \emph{Instituto Politécnico da UERJ (IPRJ-UERJ),
Nova Friburgo, Rio de Janeiro, 28625-570, Brazil.}

Eustáquio de Souza Baêta Júnior. \emph{Universidade do Estado do Rio de
Janeiro (UERJ), Faculdade de Engenharia, Programa de Pós-Graduação em
Engenharia Mecânica, Rio de Janeiro, 20550-013, Brazil.}

Letícia dos Santos Aguilera. \emph{Universidade do Estado do Rio de Janeiro
(UERJ), Instituto Politécnico (IPRJ), Programa de Pós-Graduação em Ciência
e Tecnologia de Materiais, Nova Friburgo, Rio de Janeiro, 28625-570,
Brazil.}

\textbf{Nomenclature}

\emph{BMS} = Battery Management System

\emph{BEM} = Blade Element Method

\emph{ConOps} = Concept of Operations

\emph{EPS} = Electrical Power System

\emph{FMECA} = Failure Mode, Effects, and Criticality Analysis

\emph{GPS} = Global Positioning System

\emph{IMU} = Inertial Measurement Unit

\emph{LASC} = Latin American Space Challenge

\emph{LEV} = Leading-Edge Vortex

\emph{LoRa} = Long Range wireless communication protocol

\emph{MCU} = Microcontroller Unit

\emph{PQB} = PocketQube

\emph{RBF} = Remove Before Flight

\emph{RSSI} = Received Signal Strength Indicator

\emph{SCSM} = Satellite Challenge Standards Manual

\emph{SRAB} = Sistema de Recuperação Autorrotativo Bioinspirado
(Bioinspired Autorotative Recovery System)

\emph{$\theta$} = conicity angle

\emph{$\dot{\phi}$} = rotor spin rate

\emph{$C_L$} = lift coefficient

\emph{$C_D$} = drag coefficient

\emph{$C_{D0}$} = parasitic (flat-plate) drag coefficient

\emph{$\alpha$} = local angle of attack

\emph{$\rho$} = air density

\emph{$v_0$} = vertical descent velocity

\emph{$r$} = radial position along the wing

\emph{$U$} = local resultant airspeed

\emph{$V_F$} = impact (terminal) velocity

\emph{$FS$} = safety factor

\emph{$R$} = risk score ($R = P \times S$)

\emph{$P$} = failure probability

\emph{$S$} = failure severity
